\begin{ejer}
Describe un isomorfismo $A_f\cong A[t]/(tf-1)$.
\end{ejer}
\begin{sol}
Denotemos $C\coloneq A[t]/(tf-1)$. Como $\ubar t\ubar f=1$ en $C$, el elemento $\ubar f$ tiene inverso $\ubar t$. Entonces el morfismo $A\hookrightarrow A[t]\to C$ dado por la composición de la inclusión y la proyección canónica manda $f$ al elemento $\ubar f$ invertible, por lo que induce un morfismo
\[
\Phi\colon A_f\longrightarrow C,
\qquad
\Phi\left(\frac{a}{f^n}\right)=\ubar a(\ubar f)^{-n}=\ubar a\,\ubar t^n,
\qquad \forall a\en A,\ \forall n\en\mathbb{N}.
\]
Por otro lado, el morfismo canónico $A\to A_f$ induce un morfismo (prop. universal):
\[
\theta\colon A[t]\longrightarrow A_f,
\qquad
t\longmapsto\frac{1}{f},
\qquad
a\longmapsto\frac{a}{1}
\]
Como $\theta(tf-1)=\frac{1}{f}\frac{f}{1}-1=0$, se tiene $(tf-1)\subseteq\ker\theta$, así $\theta$ induce un morfismo desde el cociente:
\[
\Psi\colon C\longrightarrow A_f,
\qquad
\ubar t\longmapsto\frac{1}{f}
\quad\text{y}\quad
\ubar a\longmapsto\frac{a}{1}
\]
Finalmente,
\[
(\Psi\circ\Phi)\left(\frac{a}{f^n}\right)
=\Psi(\ubar a\ubar t^n)
=\frac{a}{1}\frac{1}{f^n}
=\frac{a}{f^n}
\qquad\text{y}\qquad
(\Phi\circ\Psi)(\ubar a\ubar t^n)
=\Phi\left(\frac{a}{f^n}\right)
=\ubar a\ubar t^n
\]
Esto prueba que $\Phi$ es isomorfismo.
\end{sol}
